\begin{definition}[Stückweise $C^1$-Kurve\label{jordan:definition:stueckweise_c1_kurve}]
    Eine stückweise $C^1$-Kurve mit endlich vielen Ecken und Spitzen ist eine stetige injektive Abbildung $\gamma: S^1 \to \mathbb{R}^2$ mit der Eigenschaft, dass eine endliche Zerlegung $0 = t_0 < t_1 < \ldots < t_n = 1$ von $S^1$ existiert, sodass $\gamma$ auf jedem offenen Teilintervall $(t_{k}, t_{k+1})$ regulär $C^1$ ist ($\gamma'(t)$ ist stetig und $\gamma'(t) \neq 0$).
    Zusätzlich existieren an jedem $t_k$ rechtsseitige und linksseitige Tangentenvektoren die unterschiedlich sein können. 

    \begin{itemize}
        \item Eine Ecke ist ein Punkt $t_k$, an dem die rechtsseitige und linksseitige Tangente nicht kollinear sind.
        \item Eine Spitze ist ein Punkt $t_k$, an dem die rechtsseitige und linksseitige Tangente kollinear, aber entgegengesetzt gerichtet sind.
    \end{itemize}
\end{definition}